\begin{tikzpicture}[
  automaton/.style={
    draw=ciaoBlue,
    fill=ciaoPaleBlue!35,
    very thick,
    rounded corners=2mm,
    minimum width=42mm,
    minimum height=23mm,
    align=center,
    font=\sffamily
  },
  connector/.style={draw=ciaoNavy,thick},
  operation/.style={font=\ttfamily\small,align=center}
]
  \node[automaton] (view) {
    {\scriptsize\color{ciaoMuted}\textless\textless ViewAutomaton\textgreater\textgreater}\\[-1mm]
    {\bfseries SumView}
  };
  \node[automaton,right=68mm of view] (controller) {
    {\scriptsize\color{ciaoMuted}\textless\textless ControllerAutomaton\textgreater\textgreater}\\[-1mm]
    {\bfseries SumController}
  };

  % Обе половины Connection завершаются в общей точке, как при экспорте
  % в draw.io: socket required-интерфейса принимает ball provided-интерфейса.
  \coordinate (junction) at ($(view.east)!0.5!(controller.west)$);

  % Required-половина от SumView до общего стыка.
  \draw[connector] (view.east) -- ($(junction)+(-3mm,0)$);
  \draw[connector]
    ($(junction)+(0,3mm)$)
    arc[start angle=90,end angle=270,radius=3mm];

  % Provided-половина от SumController; ball касается раскрытия socket.
  \coordinate (ball-center) at ($(junction)+(2.4mm,0)$);
  \draw[connector]
    (controller.west) -- ($(ball-center)+(2.4mm,0)$);
  \draw[connector,fill=white] (ball-center) circle[radius=2.4mm];

  % Экспорт draw.io подписывает всю assembly-связь один раз у её середины.
  \node[operation,above=6mm of junction] {start};
\end{tikzpicture}
